\documentclass[a4paper,14pt]{extarticle}

\usepackage{polyglossia}
\setdefaultlanguage{russian}

\usepackage{amsmath, amssymb}
\usepackage{graphicx}
\usepackage{geometry}
\usepackage{setspace}
\usepackage{indentfirst}
\usepackage{fontspec}
\usepackage{tabularx}
\usepackage{colortbl}
\usepackage{pgfplots}

\pgfplotsset{compat=1.18}
\definecolor{lightgray}{gray}{0.9}
\setmainfont{Times New Roman}

\geometry{left=2cm,right=2cm,top=2cm,bottom=2cm}
\onehalfspacing
\setlength{\parindent}{1.25cm}

\begin{document}

\begin{titlepage}
\begin{center}
\textbf{Министерство образования Республики Беларусь}\\
\textbf{Белорусский государственный университет}

Факультет радиофизики и компьютерных технологий
\vspace{6cm}

\textbf{Лабораторная работа №13}\\
\textbf{Изучение магнитного поля соленоида и колец Гельмгольца}
\end{center}

\vfill
\begin{flushright}
Выполнила:\\
Студентка 1 курса 7 группы\\
Данильченко Мария\\
Преподаватель: Садов П. В.\\
\end{flushright}

\vfill
\begin{center}
Минск, 2026
\end{center}
\end{titlepage}

\section*{Цель работы}
Изучение распределения индукции магнитного поля на оси конечного соленоида и системы из двух катушек (колец Гельмгольца). Экспериментальная проверка теоретических законов, описывающих магнитные поля токов, и оценка однородности поля в кольцах Гельмгольца.

\section*{Метод измерений и краткая теория}
Для измерения индукции магнитного поля в данной работе применяется метод, основанный на эффекте Холла. Датчик Холла перемещается вдоль оси $z$ исследуемой системы (соленоида или колец Гельмгольца), а значения магнитной индукции $B$ фиксируются с помощью измерительного прибора.

Теоретический расчет магнитного поля кругового тока и соленоида базируется на законе Био-Савара-Лапласа:
$$d\vec{B} = \frac{\mu_0 I}{4\pi} \frac{[d\vec{l} \times \vec{r}]}{r^3}$$
где $\mu_0$ — магнитная постоянная, $I$ — сила тока, $d\vec{l}$ — элемент длины проводника, $\vec{r}$ — радиус-вектор к точке наблюдения.

Индукция магнитного поля на оси одного витка с током радиуса $R$ на расстоянии $z$ от его центра равна:
$$B_z = \frac{\mu_0 I R^2}{2(R^2 + z^2)^{3/2}}$$

\textbf{Соленоид:}
Магнитное поле конечного соленоида на его оси является суперпозицией полей отдельных витков. Для плотно намотанного соленоида формула имеет вид:
$$B = \frac{\mu_0 n I}{2} (\cos\alpha_1 + \cos\alpha_2)$$
где $n$ — число витков на единицу длины, а $\alpha_1$ и $\alpha_2$ — углы, под которыми видны торцы соленоида из точки наблюдения. В центре длинного соленоида поле максимально и считается однородным, а к краям спадает.

\textbf{Кольца Гельмгольца:}
Для получения высокооднородного магнитного поля используются кольца Гельмгольца — две одинаковые коаксиальные катушки, расположенные на расстоянии, равном их радиусу ($d = R$). Токи в катушках текут в одном направлении. Магнитная индукция в центре такой системы описывается формулой:
$$B = \left(\frac{4}{5}\right)^{3/2} \frac{\mu_0 N I}{R}$$
где $N$ — число витков в каждой катушке. В области между катушками градиент магнитного поля минимален, что создает зону однородного поля.

\section*{Результаты измерений}

\begin{center}
\small
\textbf{Таблица 3. Результаты измерений}
\begin{tabular}{|c|c|c|c|c|c|c|c|}
\hline
$z$, мм & $B$, мТл & $B_1$, мТл & $\sigma_1$, \% & $B_2$, мТл & $\sigma_2$, \% & $B_3$, мТл & $\sigma_3$, \% \\ \hline
0 & 3.52 & 3.878 & -9.23 & 3.739 & -5.86 & 3.57 & -1.4 \\ \hline
10 & 3.87 & 3.962 & -2.32 & 3.86 & 0.26 & 3.717 & 4.11 \\ \hline
20 & 4.43 & 4.196 & 5.58 & 4.156 & 6.59 & 4.09 & 8.31 \\ \hline
30 & 4.78 & 4.551 & 5.03 & 4.605 & 3.80 & 4.616 & 3.55 \\ \hline
40 & 4.85 & 4.964 & -2.3 & 5.123 & -5.33 & 5.188 & -6.52 \\ \hline
50 & 4.94 & 5.353 & -7.89 & 5.611 & -11.96 & 5.701 & -13.35 \\ \hline
60 & 4.93 & 5.684 & -13.26 & 5.982 & -17.58 & 6.068 & -18.75 \\ \hline
70 & 4.38 & 4.741 & -7.61 & 4.825 & -9.22 & 4.802 & -8.79 \\ \hline
80 & 3.86 & 3.682 & 4.84 & 3.653 & 5.67 & 3.62 & 6.63 \\ \hline
90 & 3.5 &  2.941 & 19 & 2.89 & 21.09 & 2.815 & 21.73 \\ \hline
100 & 3.05 & 2.443 & 24.85 & 2.396 & 27.31 & 2.395 & 27.34 \\ \hline
120 & 2.23 & 1.819 & 44.57 & 1.791 & 46.84 & 1.803 & 45.87 \\ \hline
130 & 2.26 & 1.609 & 40.44 & 1.586 & 42.49 & 1.599 & 41.44 \\ \hline
\end{tabular}
\end{center}

\begin{center}
\small
\textbf{Таблица 5. Результаты измерений}
\begin{tabular}{|c|c|c|c|c|c|c|}
\hline
$z$, мм & $B$, мТл & $a$, \% & $B_1$, мТл & $a_1$, \% & $B_2$, мТл & $a_2$, \% \\ \hline
0 & 3.67 & 0 & 3.68 & 0 & 3.69 & 0\\ \hline
20 & 4.13 & 12.53 & 4.15 & 12.77 & 4.16 & 12.74 \\ \hline
30 & 4.6 & 25.34 & 4.58 & 24.46 & 4.59 & 24.39 \\ \hline
40 & 6.97 & 35.42 & 4.96 & 37.78 & 4.97 & 34.69 \\ \hline
50 & 5.18 & 41.14 & 5.2 & 41.3 & 5.28 & 41.19 \\ \hline
60 & 5.15 & 40.33 & 5.17 & 40.49 & 5.18 & 40.38 \\ \hline
70 & 5.03 & 37.06 & 5.05 & 37.23 & 5.06 & 38.13 \\ \hline
80 & 4.74 & 29.16 & 4.77 & 29.62 & 4.78 & 29.54 \\ \hline
90 & 4.38 & 19.35 & 4.41 & 19.84 & 4.42 & 19.78 \\ \hline
100 & 3.89 & 6 & 4.02 & 9.24 & 4.03 & 9.21 \\ \hline
110 & 3.61 & -1.63 & 3.67 & -0.27 & 3.68 & -0.27 \\ \hline
120 & 3.29 & -10.35 & 3.37 & -8.42 & 3.38 & -8.4 \\ \hline
130 & 3.62 & -17.71 & 3.12 & -15.22 & 3.13 & -15.18 \\ \hline
140 & 2.83 & -22.89 & 2.92 & -20.65 & 2.93 & -20.59 \\ \hline
150 & 2.8 & -26.43 & 2.77 & -24.73 & 2.78 & -21.66 \\ \hline
160 & 2.64 & -28.07 & 2.65 & -28 & 2.66 & -27.92 \\ \hline
170 & 2.65 & -27.79 & 2.56 & -30.43 & 2.57 & -30.36 \\ \hline
180 & 2.73 & -25.61 & 2.49 & -32.34 & 2.5 & -32.27 \\ \hline
\end{tabular}
\end{center}

\begin{center}
\small
\textbf{Таблица 4. Результаты измерений}
\begin{tabular}{|c|c|c|c|c|c|c|}
\hline
$z$, мм & $B$, мТл & $a$, \% & $B_1$, мТл & $a_1$, \% & $B_2$, мТл & $a_2$, \% \\
\hline
0 & 4.07 & 0 & 4.07 & 0 & 4.07 & 0 \\ \hline
10 & 4.61 & 13.27 & 4.6 & 13.02 & 4.6 & 13.02 \\ \hline
20 & 5.13 & 26.04 & 5.09 & 25.06 & 5.1 & 25.13 \\ \hline
30 & 5.68 & 39.56 & 5.51 & 35.38 & 5.54 & 36.12 \\ \hline
40 & 5.98 & 46.93 & 5.84 & 43.49 & 5.89 & 44.72 \\ \hline
50 & 6.02 & 47.91 & 6.08 & 49.38 & 6.12 & 50.37 \\ \hline
60 & 6.21 & 52.58 & 6.21 & 52.58 & 6.23 & 53.07 \\ \hline
70 & 6.24 & 53.32 & 6.24 & 53.32 & 6.22 & 52.83 \\ \hline
80 & 6.17 & 51.60 & 6.17 & 51.60 & 6.10 & 49.88 \\ \hline
90 & 6.08 & 49.38 & 6.03 & 48.16 & 5.92 & 45.95 \\ \hline
100 & 5.98 & 46.93 & 5.84 & 43.49 & 5.71 & 40.29 \\ \hline
110 & 5.95 & 46.19 & 5.63 & 38.33 & 5.50 & 35.14 \\ \hline
120 & 5.87 & 44.22 & 5.43 & 33.42 & 5.31 & 30.47 \\ \hline
130 & 6.05 & 48.65 & 5.24 & 28.75 & 5.14 & 26.29 \\ \hline
140 & 6.14 & 50.86 & 5.09 & 25.06 & 5 & 22.85 \\ \hline
150 & 6.2 & 52.33 & 4.96 & 21.87 & 4.88 & 19.9 \\ \hline
160 & 6.15 & 51.11 & 4.86 & 19.41 & 4.79 & 17.69 \\ \hline
170 & 5.86 & 43.98 & 4.78 & 17.44 & 4.72 & 15.97 \\ \hline
\end{tabular}
\end{center}

\section*{Вывод}
В ходе выполнения лабораторной работы было изучено распределение магнитной индукции вдоль оси соленоида и колец Гельмгольца. 

Анализ поля соленоида (Таблица 3) показывает, что индукция магнитного поля $B$ достигает своего максимума в центральной части (при $z \approx 40 \dots 60$ мм, где $B \approx 4.9$ мТл) и плавно убывает к краям (при $z = 130$ мм поле падает до $2.26$ мТл). Сравнение с теоретическими расчетами ($B_1, B_2, B_3$) выявляет наличие относительных отклонений ($\sigma$), которые на краях возрастают. Это связано с краевыми эффектами соленоида и возможной погрешностью позиционирования датчика Холла.

Данные колец Гельмгольца (Таблица 4) наглядно демонстрируют главное свойство этой системы — создание зоны однородного поля. В диапазоне координат от $z = 60$ мм до $z = 80$ мм значения магнитной индукции остаются практически неизменными ($B \approx 6.17 \dots 6.24$ мТл), образуя характерное «плато». Измерения в Таблице 5 демонстрируют распределение поля при измененных параметрах установки, где распределение носит более неравномерный характер.

В целом, экспериментально полученные распределения качественно и количественно согласуются с теоретическими моделями, основанными на законе Био-Савара-Лапласа.

\end{document}